\section{Rúbrica}
\begin{table}[h!]
\centering
\renewcommand{\arraystretch}{1.5}
\begin{tabular}{|>{\raggedright}p{8cm}|c|c|}
\hline
\rowcolor[HTML]{D9D9D9}
\textbf{Tarea} & \textbf{Puntuación} & \textbf{Máximo} \\
\hline
\cellcolor[HTML]{E6E6E6} Presentación del informe &  & 10 \% \\
\hline
\cellcolor[HTML]{E6E6E6} Reconoce la física del JFET &  & 10 \% \\
\hline
\cellcolor[HTML]{E6E6E6} Montaje de circuitos propuestos en laboratorio e interpretación de las curvas características del JFET &  & 20 \% \\
\hline
\cellcolor[HTML]{E6E6E6} Simulación de circuitos propuestos e interpretación de las curvas características del JFET &  & 20 \% \\
\hline
\cellcolor[HTML]{E6E6E6} Entiende los límites de operación del JFET e interpreta los parámetros de la hoja de datos &  & 10 \% \\
\hline
\cellcolor[HTML]{E6E6E6} Defensa de las conclusiones &  & 30 \% \\
\hline
\rowcolor[HTML]{C0C0C0} \textbf{Total} &  & \textbf{100 \%} \\
\hline
\end{tabular}
\caption{Tareas y puntuaciones del informe sobre el JFET}
\end{table}


\newpage

\section{Introducción}
\sangria{} El presente trabajo tiene como objetivo el estudio y caracterización del transistor de efecto campo JFET, usando ensayos de los cuales se harán mediciones y previas simulaciones antes de la implementación.

\subsection{Condiciones de trabajo y ensayo}
\imagen[Temperatura ambiente en el laboratorio en ºC]{6cm}{./imagenes/temperaturaAmbiente.jpg}

